Pohled na vývoj od roku 2004, tj.\ od vstupu ČR do \ac{EU}, ukazuje, že evropská integrace sama o~sobě pokles hustoty odborů nezastavila --- naopak, otevření trhu práce přispělo k~nárůstu atypických forem zaměstnávání a~outsourcingu, které systematicky vylučují pracovníky z~dosahu stávajících odborových struktur.
Zatímco Skandinávie (DK, SE, FI) si v~tomto období udržela hustotu nad \SI{60}{\percent}, středoevropské země (CZ, SK, PL, HU) zaznamenaly souběžný pokles, přičemž CZ klesá nejrychleji.
Příznačné je, že pokles hustoty nekoreluje s~hospodářskou výkonností --- CZ v~tomto období rostla rychleji než většina srovnávaných zemí; je tedy třeba hledat institucionální příčiny, zejména absenci mechanismu automatického rozšiřování závaznosti odvětvových \ac{KS} (erga omnes), který v~AT nebo FR udržuje vysoké pokrytí i~při nízké hustotě.
Navrhovaná reforma (viz~kap.~\ref{ch:inovace}) proto cílí právě na tento mechanismus jako hlavní páku obnovy kolektivního pokrytí.

Pohled na vývoj od roku 2004, tj.\ od vstupu ČR do \ac{EU}, ukazuje, že evropská integrace sama o~sobě pokles hustoty odborů nezastavila --- naopak, otevření trhu práce přispělo k~nárůstu atypických forem zaměstnávání a~outsourcingu, které systematicky vylučují pracovníky z~dosahu stávajících odborových struktur.
Zatímco Skandinávie (DK, SE, FI) si v~tomto období udržela hustotu nad \SI{60}{\percent}, středoevropské země (CZ, SK, PL, HU) zaznamenaly souběžný pokles, přičemž CZ klesá nejrychleji.
Příznačné je, že pokles hustoty nekoreluje s~hospodářskou výkonností --- CZ v~tomto období rostla rychleji než většina srovnávaných zemí; je tedy třeba hledat institucionální příčiny, zejména absenci mechanismu automatického rozšiřování závaznosti odvětvových \ac{KS} (erga omnes), který v~AT nebo FR udržuje vysoké pokrytí i~při nízké hustotě.
Navrhovaná reforma (viz~kap.~\ref{ch:inovace}) proto cílí právě na tento mechanismus jako hlavní páku obnovy kolektivního pokrytí.

Pohled na vývoj od roku 2004, tj.\ od vstupu ČR do \ac{EU}, ukazuje, že evropská integrace sama o~sobě pokles hustoty odborů nezastavila --- naopak, otevření trhu práce přispělo k~nárůstu atypických forem zaměstnávání a~outsourcingu, které systematicky vylučují pracovníky z~dosahu stávajících odborových struktur.
Zatímco Skandinávie (DK, SE, FI) si v~tomto období udržela hustotu nad \SI{60}{\percent}, středoevropské země (CZ, SK, PL, HU) zaznamenaly souběžný pokles, přičemž CZ klesá nejrychleji.
Příznačné je, že pokles hustoty nekoreluje s~hospodářskou výkonností --- CZ v~tomto období rostla rychleji než většina srovnávaných zemí; je tedy třeba hledat institucionální příčiny, zejména absenci mechanismu automatického rozšiřování závaznosti odvětvových \ac{KS} (erga omnes), který v~AT nebo FR udržuje vysoké pokrytí i~při nízké hustotě.
Navrhovaná reforma (viz~kap.~\ref{ch:inovace}) proto cílí právě na tento mechanismus jako hlavní páku obnovy kolektivního pokrytí.

Pohled na vývoj od roku 2004, tj.\ od vstupu ČR do \ac{EU}, ukazuje, že evropská integrace sama o~sobě pokles hustoty odborů nezastavila --- naopak, otevření trhu práce přispělo k~nárůstu atypických forem zaměstnávání a~outsourcingu, které systematicky vylučují pracovníky z~dosahu stávajících odborových struktur.
Zatímco Skandinávie (DK, SE, FI) si v~tomto období udržela hustotu nad \SI{60}{\percent}, středoevropské země (CZ, SK, PL, HU) zaznamenaly souběžný pokles, přičemž CZ klesá nejrychleji.
Příznačné je, že pokles hustoty nekoreluje s~hospodářskou výkonností --- CZ v~tomto období rostla rychleji než většina srovnávaných zemí; je tedy třeba hledat institucionální příčiny, zejména absenci mechanismu automatického rozšiřování závaznosti odvětvových \ac{KS} (erga omnes), který v~AT nebo FR udržuje vysoké pokrytí i~při nízké hustotě.
Navrhovaná reforma (viz~kap.~\ref{ch:inovace}) proto cílí právě na tento mechanismus jako hlavní páku obnovy kolektivního pokrytí.

\input{texparts/python/union_density_trend_2004}



